% !TeX root = ../main.tex
% Add the above to each chapter to make compiling the PDF easier in some editors.

\chapter{Introduction}\label{chapter:introduction}

In large-scale ocean models and other grid-based scientific applications, the
computational domain is typically divided into multiple subdomains and assigned
to different Message Passing Interface (MPI) processes. To ensure the correctness
of numerical computations near subdomain boundaries, neighboring processes must
periodically exchange halo or ghost-cell data. As the degree of parallelism
increases, this neighborhood communication can introduce substantial
communication and synchronization overhead. Nonblocking MPI communication makes
it possible to perform independent computation while data is being exchanged,
and therefore offers an opportunity to hide part of the communication cost
through communication--computation overlap. Effective overlap cannot, however,
be assumed in every case. Nonblocking communication primarily permits
communication and computation to execute concurrently at the logical level,
whereas actual overlap depends on the communication progress mechanism provided
by the MPI implementation and may vary with message size and the underlying
communication mechanism~\parencite{denis2016mpi}.

This thesis investigates these effects through a benchmark-based evaluation
on the LRZ CoolMUC-4 system. The LLNL Comb halo-exchange benchmark~\parencite{combGithub} is extended
with an adjustable computation window and explicit polling points. These
extensions make it possible to examine how outstanding communication
progresses while computation is performed under the selected MPI and
transport configurations.

\section{Contributions}
\label{sec:contributions}

The contributions of this thesis are:
\begin{enumerate}
    \item A controlled experimental methodology for evaluating communication--computation overlap, implemented using a modified Comb benchmark with an adjustable synthetic computation window and optional polling during computation.
    \item A comparative evaluation of Open MPI, MPICH, and Intel MPI on LRZ CoolMUC-4, analyzing how MPI implementation, transport, message size, asynchronous progress, completion strategy, and computation duration influence the measured overlap.
    \item A web-based experiment management and visualization tool that supports experiment configuration, script generation, run management, and comparison of measured results. The tool and accompanying experimental
    materials are provided through the repository linked in
    \autoref{app:results-repository}.
\end{enumerate}

The results characterize communication--computation overlap under the evaluated benchmark, platform, and runtime configurations. They should therefore be interpreted as evidence for the tested environment rather than as direct evidence of application-level performance in a complete ocean model.

The remainder of this thesis presents the background and related work in
\autoref{chapter:background}, the benchmark methodology and experimental
setup in \autoref{chapter:methodology}, and the measurements and analysis
in \autoref{chapter:Experimental Results and Analysis}. The final chapter
summarizes the conclusions and directions for future work.
